\section{Reading One Properly}
\label{sec:ch15-reading-one-properly}

\index{resolvability!how to read one of these|(}%
Four things to do with a resolvability error, in this order. None of them
starts with the query it printed.

\index{resolvability!start at the second reason line}%
Start at the second reason line. It names the type and the field, and it names
the subgraph that declares them, which is the file to open. The path above it
is where the composer was standing rather than where the mistake is, and in
this chapter's own example it names a service that has nothing to do with
either.

\index{resolvability!ask about the key}%
Then ask that subgraph one question: can the router get in? A contributed
field is answered by sending that service a representation, so the type it
hangs off needs a key the router is allowed to enter through. Look for
\code{resolvable: false} first, because it is the argument that says no and the
one whose English reads like something else. If the key is fine, ask whether
the type has one at all, and whether the fields named in it are the ones the
other subgraph can supply.

\index{resolvability!the rule for resolvable: false}%
The rule I use for that argument is not about who owns the rows. It is about
what the type carries. A stub that declares nothing but the key exists so
another type can point at it, and nothing will ever be fetched from it, so it
takes \code{resolvable: false}. A subgraph that contributes even one field to
the type will be sent representations of it, and its key has to be enterable.
The Sessions service's \code{Speaker} is the first kind. The Search service's
\code{Session} is the first kind. The Ratings service's \code{Session} is the
second, and it is the second whether or not that service owns a single row.

\index{resolvability!fix, then compose again, then compose again}%
Third, compose again and expect another one. The check returns at the first
root field it cannot walk, so a clean run after one fix is not evidence that
there was one problem. I compose until it passes twice: once after the fix, and
once after that, having changed nothing.

\subsection{And Do Not Read a Pass as a Result}

\index{resolvability!what a pass is worth}%
The fourth thing is the one this chapter exists for.
Section~\ref{sec:ch15-the-composition-that-passes} is a graph that composes
cleanly and answers five hundred, and nothing about the input announces which
of those two you are going to get.

\index{resolvability!what a pass does establish}%
That is not an argument for ignoring the check. A failure is real every time
and is worth fixing before anything else, and merging, which is the pass that
cannot be skipped, is doing its own job throughout. What a pass establishes is
that the composer found no contradiction it was looking for. It does not
establish that a plan exists for every field, and after this chapter I would
not describe it to a colleague as though it did.

\index{resolvability!the shape that defeats the check}%
The specific shape to know about is a root field that returns an interface and
is declared shareable in more than one subgraph. If your graph has one, and a
graph with global object identification on more than one service has one
whether anybody chose it or not, then any entity reachable through that
interface can be walked into and left unresolved without a word. That is a
narrow condition and it is not an exotic one.

\index{resolvability!what to do about it}%
Two things help, and neither is a setting. Compose the pairs: a subgraph and
the one service it contributes fields to, on their own, without whatever else
is in the graph, which is a small enough input that the walk has nowhere else
to go. And test the fields rather than the composition, because
section~\ref{sec:ch15-three-fields-three-failures} is what an approved graph
does when it is wrong, and every one of those three answers would have been
caught by one request per contributed field.

\index{resolvability!the check the composer cannot run}%
Which is the uncomfortable summary. Composition reads files and never calls a
service, and chapter~\ref{ch:composition} treated that as a feature: it is why
a subgraph does not have to be running to be composed. It is also the reason
the strongest check on a graph is not composition. It is a request, sent to a
router, for a field somebody contributed to somebody else's type.
\index{resolvability!how to read one of these|)}%
